\documentclass[]{IEEEphot}

% --- Additional Packages ---
\usepackage{graphicx} % To include images (.png)
\usepackage{amsmath}  % For advanced math formulas
\usepackage{booktabs} % For professional quality tables
\usepackage{url}      % To format URLs
\usepackage{xcolor}   % For code syntax highlighting
\usepackage{listings} % To include code snippets
\usepackage{cite}     % For IEEE-style citations

% --- Configuration for listings (Python code snippets) ---
% All comments and configurations are in English as requested.
\definecolor{codegreen}{rgb}{0,0.6,0}
\definecolor{codegray}{rgb}{0.5,0.5,0.5}
\definecolor{codepurple}{rgb}{0.58,0,0.82}
\definecolor{backcolour}{rgb}{0.98,0.98,0.98}

\lstdefinestyle{mystyle}{
    backgroundcolor=\color{backcolour},   
    commentstyle=\color{codegreen},
    keywordstyle=\color{magenta},
    numberstyle=\tiny\color{codegray},
    stringstyle=\color{codepurple},
    basicstyle=\ttfamily\footnotesize,
    breakatwhitespace=false,         
    breaklines=true,                 
    captionpos=b,                    
    keepspaces=true,                 
    numbers=left,                    
    numbersep=5pt,                  
    showspaces=false,                
    showstringspaces=false,
    showtabs=false,                  
    tabsize=2,
    language=Python
}
\lstset{style=mystyle}
% --- End of listings configuration ---


% --- Journal Metadata (Placeholders) ---
\jvol{xx}
\jnum{xx}
\jmonth{November}
\pubyear{2025}

% --- Theorem Environments (from template) ---
\newtheorem{theorem}{Theorem}
\newtheorem{lemma}{Lemma}

\begin{document}

% --- Title and Author Block ---
\title{A Comparative Analysis of BiLSTM and Transformer Architectures for Dynamic Hand Gesture Recognition on the DYLEM-GRID Dataset}

\author{Marc'Antonio Lopez, Nicole Dalia Cilia}

\affil{Polytechnic University of Turin, Turin, 10129 Italy\\University of Enna "Kore", Enna, 94100 Italy}

% --- DOI Info (Placeholder) ---
\doiinfo{DOI: 10.1109/JPHOT.2025.XXXXXXX\\
1943-0655/\$25.00 \copyright 2025 IEEE}

\maketitle

% --- Received Info (Placeholder) ---
\begin{receivedinfo}%
Manuscript received October 28, 2025; revised November 7, 2025.
This work is based on the "dylem-grid" open-source project (v2.0.0, "v2025") for dynamic hand gesture recognition.
\end{receivedinfo}


% --- Abstract ---
\begin{abstract}
This paper presents a technical analysis of two state-of-the-art (SOTA) deep learning architectures for dynamic hand gesture recognition, utilizing the DYLEM-GRID dataset. The dataset, acquired via a Leap Motion Controller 2, comprises 400 samples of four distinct gestures. We implement and compare: (1) a Bidirectional Long Short-Term Memory (BiLSTM) network with a custom Attention mechanism (GestureRNN) and (2) an Encoder-only Transformer model (GestureTransformer). Both models are trained on data processed through a rigorous pipeline, including per-class outlier removal and Principal Component Analysis (PCA) retaining 95\% of variance. Results show that both optimized architectures achieve 100\% validation accuracy, effectively solving the classification task for this dataset. Notably, the GestureTransformer model (11,652 parameters) achieves this performance with 72\% fewer parameters than the GestureRNN (42,277 parameters), demonstrating significantly higher parameter efficiency.
\end{abstract}

% --- Keywords ---
\begin{IEEEkeywords}
Dynamic Hand Gesture Recognition (D-HGR), Deep Learning, Time-Series Analysis, BiLSTM, Transformer, Attention Mechanism, Leap Motion, PCA.
\end{IEEEkeywords}


% --- Start of Paper Body ---

\section{Introduction}
Dynamic Hand Gesture Recognition (D-HGR) is a fundamental field within Human-Computer Interaction (HCI), focusing on interpreting human motion over time. The primary challenge lies in capturing and modeling the complex spatial and temporal dependencies inherent in gesture sequences. While traditional methods like Hidden Markov Models (HMMs) \cite{rabiner1989} and Dynamic Time Warping (DTW) have been used, the advent of deep learning has introduced more powerful models for sequential data.

Recurrent Neural Networks (RNNs), specifically Long Short-Term Memory (LSTM) \cite{hochreiter1997} and Gated Recurrent Unit (GRU) \cite{cho2014} networks, became the standard for their ability to maintain an internal state and capture temporal context, mitigating the vanishing gradient problem. More recently, Transformer architectures \cite{vaswani2017}, originating from Natural Language Processing (NLP), have shown SOTA performance by using self-attention mechanisms to weigh the importance of different elements in a sequence, allowing for superior modeling of long-range dependencies and increased parallelization.

This project, based on the "dylem-grid" (v2025) repository, provides a direct, optimized comparison of these two leading architectures on the high-fidelity DYLEM-GRID skeletal dataset \cite{dylem2024}. We demonstrate a complete and reproducible pipeline, from raw sensor data to a highly accurate classification model. Our findings confirm that while both BiLSTM with Attention and an Encoder-only Transformer can perfectly classify the four gestures in the dataset, the Transformer architecture does so with significantly greater parameter efficiency, highlighting its advantages for this task.

\section{Related Work}
The challenge in D-HGR has always been the effective modeling of spatio-temporal dependencies. 
\subsection{Traditional Methods}
Early approaches in machine learning relied on methods like Dynamic Time Warping (DTW) for sequence alignment and Hidden Markov Models (HMMs) for probabilistic temporal modeling \cite{rabiner1989}. While effective in controlled scenarios, these methods often require extensive, manual feature engineering and struggle with the complexity of fluid, long-term gesture dynamics and high-dimensional feature spaces.

\subsection{Deep Learning Approaches}
Deep learning provided a paradigm shift by enabling end-to-end learning from sequential data.
\subsubsection{Recurrent Neural Networks}
LSTMs \cite{hochreiter1997} and their bidirectional variants (BiLSTMs) became the standard for sequential tasks. By processing sequences in both forward and backward directions, BiLSTMs can form a representation of a time step based on both past and future context. The addition of attention mechanisms \cite{bahdanau2014} further enhanced RNNs by allowing the model to learn to "focus" on the most relevant time steps in a sequence when making a classification, rather than relying on a simple average or the final hidden state.

\subsubsection{Convolutional Neural Networks}
Concurrently, 3D Convolutional Neural Networks (3D-CNNs) were employed to extract spatio-temporal features directly from video inputs (RGB or depth data), treating the time dimension as a third spatial dimension.

\subsubsection{Transformers}
The Transformer \cite{vaswani2017} model introduced an alternative approach, eschewing recurrence entirely in favor of self-attention. By calculating relationships between all parts of a sequence simultaneously, it can capture global dependencies more effectively and, crucially, can be trained in a highly parallel fashion. While dominant in NLP, its application to time-series classification, particularly in vision \cite{dosovitskiy2020} and sensor data \cite{transformer_timeseries_review}, is a growing field of research.

This paper contributes a direct benchmark of two SOTA sequential architectures (BiLSTM+Attention and Transformer) on a modern, high-resolution 3D skeletal gesture dataset \cite{dylem2024}.

\section{Dataset and Methodology}
\subsection{The DYLEM-GRID Dataset}
This work utilizes the DYLEM-GRID (Dynamic Leap Motion Gesture Recognition Indexed Dataset) \cite{dylem2024}. This dataset was created to provide a clean, indexed resource for benchmarking gesture recognition algorithms.
\begin{itemize}
    \item \textbf{Acquisition:} Data was collected using an Ultraleap Leap Motion Controller 2 sensor, capturing 3D skeletal hand data at approximately 30 Hz.
    \item \textbf{Gestures:} It includes 400 total samples of four common, cross-cultural gestures: \texttt{air\_quotes}, \texttt{finger\_wagging}, \texttt{waving}, and \texttt{zoom}.
    \item \textbf{Structure:} The dataset is pre-split into training (320 samples) and testing (80 samples) sets, ensuring reproducible evaluations.
    \item \textbf{Features:} The raw dataset (\texttt{DYLEM-GRID\_Raw}) provides 244 features per time step, detailing the kinematics (position, velocity, orientation) of the palm and 27 unique hand elements (joints and bones).
\end{itemize}

\subsection{Preprocessing Pipeline}
This project implements a rigorous, custom preprocessing pipeline (defined in \texttt{utils/data\_processing.py}) that processes the \texttt{DYLEM-GRID\_Raw} data. This pipeline is a crucial component of the methodology and is divided into four phases.

\subsubsection{Data Loading}
Samples are recursively loaded from the \texttt{DYLEM-GRID\_Raw} folders. Each sample is stored as a separate pandas DataFrame in a list, with a corresponding list of string labels (e.g., 'waving').

\subsubsection{Cleaning and Normalization}
This is the most complex phase, executed by the \texttt{data\_preprocess} function.
First, each DataFrame is cleaned of NaN values (using \texttt{bfill}) and duplicate columns.
Next, all samples are concatenated into a single global DataFrame for collective processing. Low-variance columns (where a single value occupies more than the 90\% of rows) are removed.
A robust per-class outlier removal technique is then applied, as shown in Listing \ref{lst:outlier}. This approach is more refined than global outlier removal, as it preserves the specific variations of each gesture.

\begin{lstlisting}[language=Python, caption={Per-class outlier removal logic from `data_preprocess` in `utils/data_processing.py`.}, label={lst:outlier}, captionpos=t]
for label in concatenated['_class_label'].unique():
    class_mask = concatenated['_class_label'] == label
    for column in concatenated.columns:
        if column == '_class_label': continue
        if pd.api.types.is_numeric_dtype(concatenated[column]):
            
            class_data = concatenated.loc[class_mask, column]
            mean = class_data.mean()
            std = class_data.std()
            
            if pd.isna(std) or std == 0: continue
            
            threshold = 3 * std
            # Find outliers ONLY within this class
            outliers = class_mask & \
                ((concatenated[column] - mean).abs() > threshold)
            
            # Replace outliers with the class mean
            concatenated.loc[outliers, column] = mean
\end{lstlisting}

Finally, the entire concatenated DataFrame is normalized using `MinMaxScaler` to a [0, 1] range and split back into its original list of samples.

\subsubsection{Dimensionality Reduction (PCA)}
After cleaning, the list of DataFrames is concatenated again. Principal Component Analysis (PCA) is applied, configured to retain 95\% of the explained variance, as shown in Listing \ref{lst:pca}. This reduces the feature dimensionality from ~174 to a much more manageable number (typically 22-25), maintaining essential information while reducing noise.

\begin{lstlisting}[language=Python, caption={PCA application from \texttt{apply\_pca} in \texttt{utils/data\_processing.py}.}, label={lst:pca}, captionpos=t]
concatenated = pd.concat(data, ignore_index=True)
numeric_cols = concatenated.select_dtypes(
    include=[np.number]
).columns

# Apply PCA to retain 95% of variance
variance_threshold = 0.95
pca = PCA(n_components=variance_threshold) 
transformed = pca.fit_transform(concatenated[numeric_cols])

transformed_df = pd.DataFrame(
    transformed, 
    columns=[f'PC{i+1}' for i in range(transformed.shape[1])]
)
# (Data is then split back into individual samples)
\end{lstlisting}

\subsubsection{Tensor Preparation}
The PCA-transformed data is prepared for PyTorch by the \texttt{prepare\_data} function. Text labels are converted to integers via `LabelEncoder`. Since sequences have different temporal lengths, the maximum length (`max\_len`) is identified, and all shorter sequences are post-padded with zeros (Listing \ref{lst:padding}). The final result is a single `torch.FloatTensor` of shape `(samples, max\_len, n\_features\_pca)`.

\begin{lstlisting}[language=Python, caption={Padding logic from \texttt{prepare\_data} in \texttt{utils/data\_processing.py}.}, label={lst:padding}, captionpos=t]
max_len = max(len(df) for df in data)
input_size = data[0].shape[1]

padded_data = []
for df in data:
    arr = df.values
    if len(arr) < max_len:
        # Pad with zeros at the end (post-padding)
        padding = np.zeros((max_len - len(arr), input_size))
        arr = np.vstack([arr, padding])
    padded_data.append(arr)
    
X = np.array(padded_data, dtype=np.float32)
y = np.array(encoded_labels, dtype=np.int64)

return torch.FloatTensor(X), torch.LongTensor(y), label_encoder
\end{lstlisting}


\section{Model Architectures}
We implement and compare two SOTA architectures for sequence classification.

\subsection{GestureRNN (BiLSTM + Attention)}
The GestureRNN model, defined in \texttt{bilstm\_model.py} (Listing \ref{lst:bilstm}), combines a Bidirectional LSTM with a custom attention mechanism to classify the gesture.

\begin{lstlisting}[language=Python, caption={GestureRNN (BiLSTM+Attention) class definition.}, label={lst:bilstm}, captionpos=t]
class GestureRNN(nn.Module):
    def __init__(self, input_size, hidden_size, num_layers, 
                 num_classes, dropout=0.3):
        super(GestureRNN, self).__init__()
        self.lstm = nn.LSTM(
            input_size,
            hidden_size,
            num_layers,
            dropout=dropout,
            batch_first=True,
            bidirectional=True,
        )
        # Attention mechanism
        self.attention = nn.Linear(hidden_size * 2, 1)
        # Fully connected layers
        self.fc1 = nn.Linear(hidden_size * 2, hidden_size)
        self.bn1 = nn.BatchNorm1d(hidden_size)
        self.dropout = nn.Dropout(dropout)
        self.fc2 = nn.Linear(hidden_size, num_classes)
        self.relu = nn.ReLU()

    def forward(self, x):
        # lstm_out shape: (batch, seq_len, hidden_size * 2)
        lstm_out, _ = self.lstm(x)
        
        # attention_weights shape: (batch, seq_len, 1)
        attention_scores = self.attention(lstm_out)
        attention_weights = torch.softmax(attention_scores, dim=1)
        
        # context shape: (batch, hidden_size * 2)
        context = torch.sum(attention_weights * lstm_out, dim=1)
        
        # Final classification MLP
        out = self.fc1(context)
        out = self.bn1(out)
        out = self.relu(out)
        out = self.dropout(out)
        out = self.fc2(out)
        return out
\end{lstlisting}

The forward pass proceeds as follows:
\begin{enumerate}
    \item The input tensor \texttt{(batch, max\_len, n\_features\_pca)} is fed into the \texttt{nn.LSTM} with \texttt{bidirectional=True}, producing an output of shape \texttt{(batch, max\_len, hidden\_size * 2)}.
    \item The \texttt{nn.Linear} (self.attention) layer acts as an attention mechanism, mapping the LSTM output at each time step to a single score.
    \item These scores are passed through a \texttt{torch.softmax(dim=1)} to create \texttt{attention\_weights}, representing the importance of each time step.
    \item A context vector is created by computing a weighted sum of the LSTM outputs (\texttt{torch.sum(attention\_weights * lstm\_out, dim=1)}). This compresses the variable-length sequence into a single, fixed-size vector of shape \texttt{(batch, hidden\_size * 2)}.
    \item This context vector is passed through a final MLP (with \texttt{BatchNorm1d} and \texttt{ReLU}) for classification.
\end{enumerate}

\subsection{GestureTransformer (Encoder-only)}
The GestureTransformer, defined in \texttt{transformer\_model.py} (Listing \ref{lst:transformer}), uses a standard Transformer Encoder architecture.

\begin{lstlisting}[language=Python, caption={GestureTransformer (Encoder-only) class definition.}, label={lst:transformer}, captionpos=t]
class GestureTransformer(nn.Module):
    def __init__(self, input_size, d_model=64, nhead=4, 
                 num_layers=2, dim_feedforward=128, 
                 num_classes=4, dropout=0.1):
        super(GestureTransformer, self).__init__()
        self.input_projection = nn.Linear(input_size, d_model)
        self.pos_encoder = PositionalEncoding(d_model, dropout)
        
        encoder_layer = nn.TransformerEncoderLayer(
            d_model=d_model,
            nhead=nhead,
            dim_feedforward=dim_feedforward,
            dropout=dropout,
            batch_first=True,
        )
        self.transformer_encoder = nn.TransformerEncoder(
            encoder_layer, num_layers=num_layers
        )
        
        self.classifier = nn.Sequential(
            nn.Linear(d_model, dim_feedforward),
            nn.ReLU(),
            nn.Dropout(dropout),
            nn.Linear(dim_feedforward, num_classes),
        )

    def forward(self, x):
        # x shape: (batch, seq_len, input_size)
        
        # 1. Project to model dimension
        x = self.input_projection(x) # (batch, seq_len, d_model)
        
        # 2. Add positional encoding
        x = self.pos_encoder(x)
        
        # 3. Process with Self-Attention
        x = self.transformer_encoder(x) # (batch, seq, d_model)
        
        # 4. Global Average Pooling (temporal compression)
        x = torch.mean(x, dim=1) # (batch, d_model)
        
        # 5. Classify
        x = self.classifier(x) # (batch, num_classes)
        return x
\end{lstlisting}

The forward pass proceeds as follows:
\begin{enumerate}
    \item The input tensor \texttt{(batch, max\_len, n\_features\_pca)} is projected to the model's internal dimension (\texttt{d\_model}) by an initial \texttt{nn.Linear} layer.
    \item A sinusoidal \texttt{PositionalEncoding} is added to the projected input to provide temporal information. The encoding follows the original formulas:
    \begin{equation}
    PE_{(pos, 2i)} = \sin(pos / 10000^{2i / d_{\text{model}}})
    \end{equation}
    \begin{equation}
    PE_{(pos, 2i+1)} = \cos(pos / 10000^{2i / d_{\text{model}}})
    \end{equation}
    \item The resulting sequence \texttt{(batch, max\_len, d\_model)} is processed by a stack of \texttt{nn.TransformerEncoderLayer} blocks, which apply multi-head self-attention.
    \item The temporal dimension is compressed using Global Average Pooling (\texttt{torch.mean(x, dim=1)}), which averages the feature vectors across all time steps.
    \item The resulting \texttt{(batch, d\_model)} vector is passed through a final MLP for classification.
\end{enumerate}


\section{Experiments and Results}

\subsection{Training and Optimization}
The training procedure followed a two-phase process for each model.

First, a hyperparameter optimization (HPO) search was conducted using Optuna, as detailed in the \texttt{scripts/optimization/} files. The search ran for 100 trials, using a \texttt{MedianPruner} to terminate unpromising runs early. The objective was to maximize validation accuracy.

Second, the best-performing hyperparameters from the HPO phase were hardcoded into the final training scripts (\texttt{scripts/train/}). Both models were trained for a maximum of 50 epochs using a \texttt{train\_test\_split} with \texttt{random\_state=44} and \texttt{stratify=y}. Early stopping with \texttt{patience=15} was used to halt training if validation accuracy did not improve. The BiLSTM was trained with NAdam (lr=0.00195...) and the Transformer with Adam (lr=0.00108...). Both used \texttt{nn.CrossEntropyLoss()}.

\subsection{Results and Discussion}
Both models, when trained with their optimized hyperparameters, were able to perfectly solve the classification task for the validation set, achieving 100\% accuracy.

\begin{table*}[t]
\centering
\caption{Model Architecture and Performance Comparison}
\label{tab:results}
\begin{tabular}{@{}lcccc@{}}
\toprule
\textbf{Model Architecture} & \textbf{Validation Accuracy} & \textbf{Total Parameters} & \textbf{Key Hyperparameters} & \textbf{Optimizer} \\
\midrule
GestureRNN (BiLSTM+Attn) & 100.00\% & 42,277 & $\begin{array}{@{}l@{}}\text{hidden\_size} = 64 \\ \text{num\_layers} = 2 \\ \text{dropout} = 0.156\end{array}$ & NAdam \\
\addlinespace
GestureTransformer & 100.00\% & 11,652 & $\begin{array}{@{}l@{}}\text{d\_model} = 64, \text{nhead} = 4 \\ \text{num\_layers} = 1 \\ \text{dropout} = 0.494\end{array}$ & Adam \\
\bottomrule
\end{tabular}
\end{table*}

\begin{figure*}[t]
\centering
\begin{minipage}{0.48\textwidth}
    \centering
    \includegraphics[width=\textwidth]{plots/bilstm_training_history.png}
    \caption{BiLSTM (GestureRNN) training and validation accuracy/loss curves. The model reached 100\% validation accuracy at Epoch 7.}
    \label{fig:bilstm_history}
\end{minipage}\hfill
\begin{minipage}{0.48\textwidth}
    \centering
    \includegraphics[width=\textwidth]{plots/transformer_training_history.png}
    \caption{Transformer (GestureTransformer) training and validation accuracy/loss curves. The model reached 100\% validation accuracy at Epoch 20.}
    \label{fig:transformer_history}
\end{minipage}
\end{figure*}

\begin{figure*}[t]
\centering
\begin{minipage}{0.48\textwidth}
    \centering
    \includegraphics[width=\textwidth]{plots/bilstm_confusion_matrix.png}
    \caption{BiLSTM (GestureRNN) confusion matrix and per-class metrics for the best model (Epoch 7), showing perfect classification.}
    \label{fig:bilstm_cm}
\end{minipage}\hfill
\begin{minipage}{0.48\textwidth}
    \centering
    \includegraphics[width=\textwidth]{plots/transformer_confusion_matrix.png}
    \caption{Transformer (GestureTransformer) confusion matrix and per-class metrics for the best model (Epoch 20), showing perfect classification.}
    \label{fig:transformer_cm}
\end{minipage}
\end{figure*}

The key results are summarized in Table \ref{tab:results}. As shown in Fig. \ref{fig:bilstm\_history} and Fig. \ref{fig:transformer\_history}, both models converge quickly to perfect accuracy on the validation set. The BiLSTM model (Fig. \ref{fig:bilstm\_history}) achieved this faster, at epoch 7, while the Transformer (Fig. \ref{fig:transformer\_history}) reached it at epoch 20.

The confusion matrices (Fig. \ref{fig:bilstm\_cm} and Fig. \ref{fig:transformer\_cm}) confirm that the best-performing checkpoints for both models resulted in zero misclassifications on the 80-sample validation set.

The most significant finding is the difference in parameter efficiency. The GestureTransformer model achieves the same SOTA performance with only **11,652** parameters, a **72.3\% reduction** compared to the 42,277 parameters of the GestureRNN. This suggests that for this type of high-dimensional, sequential skeletal data, the self-attention mechanism is a more parameter-efficient method for capturing temporal dependencies than a recurrent architecture.

\section{Conclusion}
This paper presented a complete, robust, and SOTA pipeline for dynamic hand gesture recognition on the DYLEM-GRID dataset, comparing two leading architectures: BiLSTM with Attention (GestureRNN) and an Encoder-only Transformer (GestureTransformer).

Our findings demonstrate that both architectures, when properly optimized, can achieve 100\% validation accuracy, effectively solving the classification task for the four gestures in this dataset. The comparative analysis, however, reveals a clear advantage in efficiency for the Transformer. By achieving identical performance with 72\% fewer parameters, the GestureTransformer proves to be a more efficient and scalable solution, likely due to the efficacy of its self-attention mechanism in modeling spatio-temporal dependencies from skeletal data.

Future work, as suggested in the project's documentation, could include applying data augmentation techniques (e.g., time-warping) to test generalization, using Explainability (XAI) methods to interpret the models' decisions, and benchmarking against other time-series architectures such as Temporal Convolutional Networks (TCN).

% --- Bibliography ---
% Using \begin{thebibliography} for a self-contained example.
% A .bib file (e.g., using BibTeX) is recommended for larger papers.
\begin{thebibliography}{9}

\bibitem{dylem2024}
M. Sorce, M. A. Lopez, G. Trovato, and N. D. Cilia, "DYLEM-GRID - Dynamic Leap Motion Gesture Recognition Indexed Dataset," \emph{[Placeholder Journal/Conference]}, 2024.

\bibitem{vaswani2017}
A. Vaswani, N. Shazeer, N. Parmar, J. Uszkoreit, L. Jones, A. N. Gomez, L. Kaiser, and I. Polosukhin, "Attention is all you need," in \emph{Advances in Neural Information Processing Systems (NIPS)}, 2017, pp. 5998--6008.

\bibitem{hochreiter1997}
S. Hochreiter and J. Schmidhuber, "Long short-term memory," \emph{Neural Computation}, vol. 9, no. 8, pp. 1735--1780, 1997.

\bibitem{rabiner1989}
L. R. Rabiner, "A tutorial on hidden Markov models and selected applications in speech recognition," \emph{Proceedings of the IEEE}, vol. 77, no. 2, pp. 257--286, 1989.

\bibitem{cho2014}
K. Cho, B. Van Merriënboer, C. Gulcehre, D. Bahdanau, F. Bougares, H. Schwenk, and Y. Bengio, "Learning phrase representations using RNN encoder-decoder for statistical machine translation," in \emph{Proc. EMNLP}, 2014.

\bibitem{bahdanau2014}
D. Bahdanau, K. Cho, and Y. Bengio, "Neural machine translation by jointly learning to align and translate," in \emph{Proc. ICLR}, 2015.

\bibitem{dosovitskiy2020}
A. Dosovitskiy, L. Beyer, A. Kolesnikov, D. Weissenborn, X. Zhai, T. Unterthiner, M. Dehghani, M. Minderer, G. Heigold, S. Gelly, J. Uszkoreit, and N. Houlsby, "An image is worth 16x16 words: Transformers for image recognition at scale," in \emph{Proc. ICLR}, 2021.

\bibitem{transformer_timeseries_review}
A. A. Zerveas, S. Jayaraman, D. Patel, A. Bhamidipaty, and C. Eickhoff, "A Transformer-based framework for multivariate time series representation learning," in \emph{Proc. 38th Int. Conf. Machine Learning (ICML)}, 2021, pp. 12948--12959.

\end{thebibliography}

\end{document}